% GENERATED FILE. Edit canonical lessons, then run npm run generate:books.
% canonical-source-hash: fnv1a64:4202c1371705824f
% canonical-lessons: SA-C454-irsyati, SA-C454-jugupsate, SA-C454-dvesah, SA-C454-vairam, SA-C454-sprha

\chapter{Envy and Dislike}
\label{ch:sa-envy-and-dislike}
\hlchaptermodality{454}

\begin{chapteropening}
\textbf{By the end of this chapter:} \emph{I can say whom someone envies or dislikes.}

Complete the last lesson of chapter 454: I can say whom someone envies or dislikes.
\end{chapteropening}

\section[\texorpdfstring{\sk{ईर्ष्यति} (īrṣyati)}{ईर्ष्यति (īrṣyati)}]{\sk{ईर्ष्यति} (īrṣyati) — envies}

\label{lesson:SA-C454-irsyati}



\begin{quote}
Before the new one: say the Sanskrit for a dear friend, then the Sanskrit for sends word.
\end{quote}

\subsection*{You'll want to know: \sk{ईर्ष्यति}}
\textbf{\sk{ईर्ष्यति}} — \emph{īrṣyati} — \textquotedblleft{}envies\textquotedblright{}.

\begin{quote}
\sk{सः} \sk{धनिकाय} \sk{मित्राय} \sk{ईर्ष्यति।}
\end{quote}

\begin{quote}
\emph{saḥ dhanikāya mitrāya īrṣyati.}
\end{quote}

\begin{quote}
He envies his rich friend.
\end{quote}

Past (he …): \sk{सः} \sk{मित्राय} \sk{ईर्ष्यितवान्}, \emph{saḥ mitrāya īrṣyitavān}.

\begin{cousinweb}
\textbf{\sk{ईर्ष्यति}} (\emph{īrṣyati}) comes from the root √īrṣy, \textquotedblleft{}to envy, to feel impatient at another's prosperity\textquotedblright{}, which takes the dative; MW gives no further derivation, so its origin is uncertain (MW).
\end{cousinweb}

\begin{grammarlens}[title={What you've built}]
One more word for describing what you lived and felt.
\end{grammarlens}

\subsection*{Your turn}
\begin{itemize}
\raggedright
  \item \emph{Say it:} \emph{īrṣyati}
  \item \emph{Say it:} \emph{īrṣyati}, once more
  \item \emph{Say it:} \emph{saḥ dhanikāya mitrāya īrṣyati.} — He envies his rich friend.
  \item \emph{Recall:} say the Sanskrit for a dear friend, then the Sanskrit for sends word, then say \emph{īrṣyati} again
\end{itemize}

\begin{tcolorbox}[breakable,colback=teal!4,colframe=teal!35!black,title={Before you move on}]
What does \sk{ईर्ष्यति} mean? (\textquotedblleft{}Envies\textquotedblright{}.) Say it once more.
\end{tcolorbox}
\section[\texorpdfstring{\sk{जुगुप्सते} (jugupsate)}{जुगुप्सते (jugupsate)}]{\sk{जुगुप्सते} (jugupsate) — detests}

\label{lesson:SA-C454-jugupsate}



\begin{quote}
Before the new one: say the Sanskrit for sends word, then the Sanskrit for envies.
\end{quote}

\subsection*{You'll want to know: \sk{जुगुप्सते}}
\textbf{\sk{जुगुप्सते}} — \emph{jugupsate} — \textquotedblleft{}detests\textquotedblright{}.

\begin{quote}
\sk{मीरा} \sk{अनुचितं} \sk{कार्यं} \sk{जुगुप्सते।}
\end{quote}

\begin{quote}
\emph{mīrā anucitaṃ kāryaṃ jugupsate.}
\end{quote}

\begin{quote}
Mira detests wrong deeds.
\end{quote}

Past (she …): \sk{सा} \sk{अनुचितं} \sk{कार्यं} \sk{जुगुप्सितवती}, \emph{sā anucitaṃ kāryaṃ jugupsitavatī}.

\begin{cousinweb}
\textbf{\sk{जुगुप्सते}} (\emph{jugupsate}) is built on the desiderative of the root √gup, \textquotedblleft{}to guard\textquotedblright{}: wanting to keep oneself from a thing; MW's ppp \emph{jugupsita} is \textquotedblleft{}disliked, detested\textquotedblright{} (MW).
\end{cousinweb}

\begin{grammarlens}[title={What you've built}]
One more word for describing what you lived and felt.
\end{grammarlens}

\subsection*{Your turn}
\begin{itemize}
\raggedright
  \item \emph{Say it:} \emph{jugupsate}
  \item \emph{Say it:} \emph{jugupsate}, once more
  \item \emph{Say it:} \emph{mīrā anucitaṃ kāryaṃ jugupsate.} — Mira detests wrong deeds.
  \item \emph{Recall:} say the Sanskrit for sends word, then the Sanskrit for envies, then say \emph{jugupsate} again
\end{itemize}

\begin{tcolorbox}[breakable,colback=teal!4,colframe=teal!35!black,title={Before you move on}]
What does \sk{जुगुप्सते} mean? (\textquotedblleft{}Detests\textquotedblright{}.) Say it once more.
\end{tcolorbox}
\section[\texorpdfstring{\sk{द्वेषः} (dveṣaḥ)}{द्वेषः (dveṣaḥ)}]{\sk{द्वेषः} (dveṣaḥ) — hatred}

\label{lesson:SA-C454-dvesah}



\begin{quote}
Before the new one: say the Sanskrit for envies, then the Sanskrit for detests.
\end{quote}

\subsection*{You'll want to know: \sk{द्वेषः}}
\textbf{\sk{द्वेषः}} — \emph{dveṣaḥ} — \textquotedblleft{}hatred\textquotedblright{}.

\begin{quote}
\sk{द्वेषः} \sk{मित्रम्} \sk{अपि} \sk{शत्रुं} \sk{करोति।}
\end{quote}

\begin{quote}
\emph{dveṣaḥ mitram api śatruṃ karoti.}
\end{quote}

\begin{quote}
Hatred turns even a friend into an enemy.
\end{quote}

A masculine noun.

\begin{cousinweb}
\textbf{\sk{द्वेषः}} (\emph{dveṣaḥ}) comes from the root √dviṣ, \textquotedblleft{}to hate, to dislike\textquotedblright{}; MW gives \textquotedblleft{}hatred, dislike, repugnance, enmity\textquotedblright{} (MW).
\end{cousinweb}

\begin{grammarlens}[title={What you've built}]
One more word for describing what you lived and felt.
\end{grammarlens}

\subsection*{Your turn}
\begin{itemize}
\raggedright
  \item \emph{Say it:} \emph{dveṣaḥ}
  \item \emph{Say it:} \emph{dveṣaḥ}, once more
  \item \emph{Say it:} \emph{dveṣaḥ mitram api śatruṃ karoti.} — Hatred turns even a friend into an enemy.
  \item \emph{Recall:} say the Sanskrit for envies, then the Sanskrit for detests, then say \emph{dveṣaḥ} again
\end{itemize}

\begin{tcolorbox}[breakable,colback=teal!4,colframe=teal!35!black,title={Before you move on}]
What does \sk{द्वेषः} mean? (\textquotedblleft{}Hatred\textquotedblright{}.) Say it once more.
\end{tcolorbox}
\section[\texorpdfstring{\sk{वैरम्} (vairam)}{वैरम् (vairam)}]{\sk{वैरम्} (vairam) — enmity}

\label{lesson:SA-C454-vairam}



\begin{quote}
Before the new one: say the Sanskrit for detests, then the Sanskrit for hatred.
\end{quote}

\subsection*{You'll want to know: \sk{वैरम्}}
\textbf{\sk{वैरम्}} — \emph{vairam} — \textquotedblleft{}enmity\textquotedblright{}.

\begin{quote}
\sk{सिंहस्य} \sk{वृकस्य} \sk{च} \sk{वैरम्} \sk{अस्ति।}
\end{quote}

\begin{quote}
\emph{siṃhasya vṛkasya ca vairam asti.}
\end{quote}

\begin{quote}
There is enmity between the lion and the wolf.
\end{quote}

A neuter noun.

\begin{cousinweb}
\textbf{\sk{वैरम्}} (\emph{vairam}) is built on \emph{vīra}, \textquotedblleft{}a hero, a fighting man\textquotedblright{}, by lengthening its first vowel; MW gives \textquotedblleft{}enmity, hostility, a feud\textquotedblright{} (MW).
\end{cousinweb}

\begin{grammarlens}[title={What you've built}]
One more word for describing what you lived and felt.
\end{grammarlens}

\subsection*{Your turn}
\begin{itemize}
\raggedright
  \item \emph{Say it:} \emph{vairam}
  \item \emph{Say it:} \emph{vairam}, once more
  \item \emph{Say it:} \emph{siṃhasya vṛkasya ca vairam asti.} — There is enmity between the lion and the wolf.
  \item \emph{Recall:} say the Sanskrit for detests, then the Sanskrit for hatred, then say \emph{vairam} again
\end{itemize}

\begin{tcolorbox}[breakable,colback=teal!4,colframe=teal!35!black,title={Before you move on}]
What does \sk{वैरम्} mean? (\textquotedblleft{}Enmity\textquotedblright{}.) Say it once more.
\end{tcolorbox}
\section[\texorpdfstring{\sk{स्पृहा} (spṛhā)}{स्पृहा (spṛhā)}]{\sk{स्पृहा} (spṛhā) — covetous longing}

\label{lesson:SA-C454-sprha}



\begin{quote}
Before the new one: say the Sanskrit for hatred, then the Sanskrit for enmity.
\end{quote}

\subsection*{You'll want to know: \sk{स्पृहा}}
\textbf{\sk{स्पृहा}} — \emph{spṛhā} — \textquotedblleft{}covetous longing\textquotedblright{}.

\begin{quote}
\sk{धनस्य} \sk{स्पृहा} \sk{सन्तोषं} \sk{हरति।}
\end{quote}

\begin{quote}
\emph{dhanasya spṛhā santoṣaṃ harati.}
\end{quote}

\begin{quote}
Longing for money takes contentment away.
\end{quote}

A feminine noun.

\begin{cousinweb}
\textbf{\sk{स्पृहा}} (\emph{spṛhā}) comes from the root √spṛh, \textquotedblleft{}to desire eagerly\textquotedblright{}; MW gives \textquotedblleft{}eager desire, covetousness, envy, longing for\textquotedblright{} (MW).
\end{cousinweb}

\begin{grammarlens}[title={What you've built}]
One more word for describing what you lived and felt.
\end{grammarlens}

\subsection*{Your turn}
\begin{itemize}
\raggedright
  \item \emph{Say it:} \emph{spṛhā}
  \item \emph{Say it:} \emph{spṛhā}, once more
  \item \emph{Say it:} \emph{dhanasya spṛhā santoṣaṃ harati.} — Longing for money takes contentment away.
  \item \emph{Recall:} say the Sanskrit for hatred, then the Sanskrit for enmity, then say \emph{spṛhā} again
\end{itemize}

\begin{tcolorbox}[breakable,colback=teal!4,colframe=teal!35!black,title={Before you move on}]
What does \sk{स्पृहा} mean? (\textquotedblleft{}Covetous longing\textquotedblright{}.) Say it once more.
\end{tcolorbox}
